\documentclass[10pt,twocolumn,a4paper]{article}
\usepackage[utf8]{inputenc}
\usepackage[T1]{fontenc}
\usepackage{newtxtext,newtxmath}\renewcommand{\ttdefault}{lmtt}
\usepackage[margin=17mm,top=20mm,bottom=20mm,columnsep=7mm]{geometry}
\usepackage[protrusion=true,expansion=true]{microtype}
\usepackage{booktabs,graphicx,xcolor,enumitem,caption,titlesec,fancyhdr,lastpage,array,tabularx,amsmath}
\usepackage{tikz}
\usetikzlibrary{arrows.meta,positioning,calc,fit}
\usepackage[hyphens]{url}\Urlmuskip=0mu plus 2mu
\usepackage[authoryear,round,sort&compress]{natbib}
\usepackage[colorlinks=true,linkcolor=sigred,citecolor=coolblue,urlcolor=coolblue,pdfborder={0 0 0},
  pdftitle={From Gauge to Shelter: A Public-Data Audit of Flood and Haze Decision Support in Chiang Mai},pdfauthor={Poon Thiengburanathum; Non Arkaraprasertkul},
  pdfkeywords={smart city; digital government; Thailand; early warning; open data; digital twin; inter-agency coordination}]{hyperref}

\definecolor{ink}{HTML}{1B1F23}\definecolor{sigred}{HTML}{B4231A}\definecolor{coolblue}{HTML}{1E4E79}
\definecolor{mid}{HTML}{6B7280}\definecolor{rule}{HTML}{D4D7DC}
\renewcommand{\topfraction}{0.9}\renewcommand{\bottomfraction}{0.6}\renewcommand{\textfraction}{0.08}\renewcommand{\floatpagefraction}{0.75}\renewcommand{\dblfloatpagefraction}{0.75}\renewcommand{\dbltopfraction}{0.9}
\setlength{\parindent}{0.9em}\setlength{\parskip}{0pt}\linespread{1.04}
\setlist{nosep,leftmargin=1.3em,itemsep=1.5pt}
\captionsetup{font={small},labelfont={bf,color=sigred},labelsep=period,justification=raggedright,singlelinecheck=false,skip=3pt}
\titleformat{\section}{\large\bfseries\color{ink}}{\thesection}{0.6em}{}
\titleformat{\subsection}{\normalsize\bfseries\color{ink}}{\thesubsection}{0.6em}{}
\titlespacing*{\section}{0pt}{11pt}{4pt}\titlespacing*{\subsection}{0pt}{7pt}{2pt}
\pagestyle{fancy}\fancyhf{}\renewcommand{\headrulewidth}{0.3pt}\renewcommand{\headrule}{\hbox to\headwidth{\color{rule}\leaders\hrule height \headrulewidth\hfill}}
\fancyhead[L]{\small\color{mid}Preprint --- not peer reviewed}\fancyhead[R]{\small\color{mid}Thiengburanathum \& Arkaraprasertkul}
\setlength{\headheight}{15pt}
\fancyfoot[C]{\small\color{mid}\thepage\,/\,\pageref{LastPage}}
\fancypagestyle{first}{\fancyhf{}\renewcommand{\headrulewidth}{0pt}\fancyhead[L]{\small\color{mid}Preprint --- not peer reviewed \quad|\quad v1.1 (4 Nov 2026)}\fancyfoot[C]{\small\color{mid}\thepage\,/\,\pageref{LastPage}}}
\newcommand{\kw}[1]{\textbf{#1}}
\newcommand{\hyp}[1]{\textbf{\color{sigred}#1}}
\renewcommand{\arraystretch}{1.12}

\begin{document}
\twocolumn[{%
\thispagestyle{first}
\begin{center}
{\color{mid}\small\textsc{Preprint \,·\, Public-policy and urban-systems research}}\\[7pt]
{\LARGE\bfseries\color{ink} From Gauge to Shelter: A Public-Data Audit of Flood and Haze Decision Support in Chiang Mai\par}
\vspace{7pt}
{\large Poon Thiengburanathum$^{1,2}$ \quad and \quad Non Arkaraprasertkul$^{3}$}\\[3pt]
{\small $^{1}$Faculty of Engineering, Chiang Mai University, Thailand\\
$^{2}$Research and Innovation Acceleration Agency for Competitiveness and Area Development (RCAD), Thailand\\
$^{3}$Digital Economy Promotion Agency of Thailand (depa)}\\[2pt]
{\small\color{mid}Correspondence: \href{https://cnx.nonarkara.org/cnx/research}{cnx.nonarkara.org/cnx/research} \,·\, Preprint of 4 November 2026 \,·\, Version v1.1 (4 Nov 2026)}
\end{center}
\vspace{2pt}
\begin{center}\begin{minipage}{0.92\textwidth}\small
\textbf{Abstract.}\; Public measurements can support disaster decisions only when warnings, local assets and accountable response are connected. We examine four explanations: weak inter-agency coordination, procurement that buys files rather than services, data that records rather than responds, and no end-to-end monitoring-to-push chain. In one province we describe what the open-data estate publishes (311 datasets, 2{,}869 resources, 45 publishers), audit two hazard chains link by link, and report the operating record of CNX, a public-data operations room adapted to close the chain without inter-agency agreements. 41\% of resources could not be parsed and 6 datasets carry a location; title-and-tag searches found none about drainage, evacuation routes or procurement in the two archived catalogues. None of the 177 official shelter rows carries a coordinate, and 50\% of rows reporting toilets exceed the medium-term reference of 20 people per toilet at listed capacity. CNX drew on 29 provider labels (including composite and internally derived sources) within seven calendar weeks; in one archived snapshot, 24 of 25 feeds with a declared freshness limit were inside it; official meteorological warnings lead its board. The prototype demonstrates integration feasibility, not warning delivery, acknowledged action, safe evacuation or institutional reform. The audit is descriptive and the system is evaluated for feasibility only: we show that the gaps are present and measurable, not that they cause the stalled programmes.\\[4pt]
\textbf{Keywords:} smart city; digital government; Thailand; early warning; open data; digital twin; inter-agency coordination; procurement.
\end{minipage}\end{center}
\vspace{4pt}\hrule height 0.4pt\vspace{8pt}
}]

\section{A gauge, a river, and a calm screen}
A rain gauge and a river gauge answer different questions. The CNX test suite includes a Chom Thong rainfall observation of 94.5\,mm in twenty-four hours, alongside the problem of a river-only view that could miss heavy rain in an upstream catchment (\path{rain-core.test.ts}). This fixture illustrates a design risk; it does not establish the date or operational consequences of a historical incident. We use it to ask what a public operations room must connect before a measurement can support action.

The availability of sensors alone does not establish that connection. Thailand has rain gauges, satellites, a national meteorological service that publishes its warnings as machine-readable CAP messages \citep{tmd2026cap}, and a provincial office that has written down 177 places where flood evacuees may sleep. It is a failure of \emph{the chain}: the steps between a measurement and a person who moves. Each link is published, if at all, by a different body, in a different form.

Thailand has promoted smart cities and digital government for a decade. The United Nations e-government index places it 52nd of 193 in 2024, up from 55th in 2022, while its e-participation rank fell from 18th to 42nd over the same two years \citep{undesa2024egov}. Studies find a smart-city strategy with no province-level implementation plan \citep{naprathansuk2017} and, in Phuket, progress in only some dimensions with budget the main obstacle \citep{jamnong2022phuket}. The pattern is familiar from critiques of smart cities elsewhere \citep{hollands2008,soderstrom2014,shelton2015}.

\paragraph{Research question.} \emph{Can a working chain from measurement to warning to shelter be assembled from public sources alone, without waiting for inter-agency agreement --- and what do the points where it cannot be assembled show about why the chain stays open?}

\paragraph{Four candidate gaps.} Practitioners offer four explanations. We treat each as a hypothesis, state what we could and could not measure for it, and report descriptively:
\begin{itemize}
\item \hyp{H1 Coordination.} Agencies publish in isolation; nothing joins their records.
\item \hyp{H2 Procurement.} Technology is bought as files and projects, not as open, maintained services.
\item \hyp{H3 Responsiveness.} Published data records the past; it does not drive a response.
\item \hyp{H4 End-to-end.} No monitoring-to-push chain exists, e.g.\ rain to drainage, haze to shelter.
\end{itemize}
\paragraph{What this paper is, and is not.} It is an audit of what one province's open data \emph{publishes}, plus a feasibility study of a system built to close the chain. It does not observe procurement, agency behaviour or warning outcomes directly, and it does not test the hypotheses; Section~\ref{sec:limits} says what would.
\paragraph{Contributions.} (i) A descriptive audit of one province's open-data estate (311 datasets, 2{,}869 resources, 45 publishers) against the four gaps. (ii) A link-by-link audit of two hazard chains. (iii) A working system and its operating record, with an account of what it could not bypass. (iv) Open data and code to repeat every number.

\section{Background}
\paragraph{Smart cities and the delivery gap.}
Critics describe the smart city as an idea that sells faster than it ships \citep{hollands2008,townsend2013smart,greenfield2013against}: corporate storytelling \citep{soderstrom2014,wiig2015}, top-down programmes that skip the citizen \citep{datta2015,cardullo2019}, and a gap between rhetoric and the ``actually existing'' city \citep{shelton2015}. The label travels as a global discourse \citep{joss2019}. For Thailand, recent work identifies the core data sets that smart-city governance needs across 36 cities \citep{ruangwicha2026}.

\paragraph{Coordination and procurement.}
Joined-up government is easy to want and hard to build \citep{pollitt2003,bardach1998}; the literature lists conditions for collaboration \citep{bryson2006,ansell2008}, network forms \citep{provan2008} and whole-of-government reform \citep{christensen2007}. Public IT fails often \citep{goldfinch2007,anthopoulos2016}; project cost and design-reality gaps are recurring causes \citep{flyvbjerg2013,heeks2006}; demand-side procurement can drive innovation \citep{edler2007}; and digital-service teams offer an alternative delivery model \citep{mergel2019,pahlka2023}. Thailand's procurement act was enacted in 2017 \citep{thailand2017procurementact}; the Digital Government Act concerns digital administration and service delivery \citep{thailand2019digitalact}.

\paragraph{Responsive data.}
Open data has documented benefits, barriers and myths \citep{janssen2012open}; its quality can be measured along dimensions such as currency and accessibility \citep{vetro2016,zaveri2016}, and the FAIR principles say what a dataset must offer a machine \citep{wilkinson2016}. A city digital twin adds the demand that a model stay synchronised with what it models \citep{batty2018,bolton2018gemini}.

\paragraph{Early warning.}
An early-warning system has four elements: knowing the risk, detecting and forecasting, communicating, and preparing to respond \citep{wmo2018mhews}; the field has long treated it as a people-centred system rather than a technology \citep{basher2006,sorensen2000}. The Sendai Framework and Early Warnings for All make it a global target \citep{unga2015sendai,wmo2023ew4all}, and the Common Alerting Protocol is the open format that carries a warning between systems \citep{oasis2010cap}. Thailand's 2011 floods caused THB\,1.43 trillion in damage and losses and affected over 13 million people; the World Bank's assessment recommended early warning built on high-quality real-time data \citep{worldbank2012thaiflood}.

\paragraph{Haze.}
In Chiang Mai the other seasonal hazard is smoke \citep{chantara2012,inlaung2022}. Satellite fire detection runs at 375\,m \citep{schroeder2014}; the World Health Organization's 2021 guideline for 24-hour PM$_{2.5}$ is 15\,\textmu g/m$^3$, with interim targets down from 75 \citep{who2021aqg}; a regional treaty has required monitoring since 2002 \citep{asean2002haze}.

\section{Method}
\paragraph{Design.} Two parts. \emph{Audit}: we describe what published records show for each gap. \emph{Build}: we construct a system on public sources and record what it takes and what it cannot do. Table~\ref{tab:hyp} states, for each hypothesis, what we measured and what we could not.
\paragraph{Terms.} The \emph{estate} is the set of datasets returned by a search of the national open-data portal for the Thai word for Chiang Mai. The \emph{catalogue} is the list of resources on the provincial portal. The \emph{board} is CNX's web screen. \emph{Baked} means a dataset was fetched and frozen on a stated date.

\begin{table*}[!t]
\centering\small
\caption{Four gaps: what we measured, what we could not, and what the audit shows (Section~\ref{sec:results}). The audit is descriptive; it does not test the hypotheses.}\label{tab:hyp}
\begin{tabularx}{\textwidth}{@{}p{2.2cm}XX@{}}
\toprule
\textbf{Gap} & \textbf{Measured / not measured} & \textbf{What the audit shows}\\\midrule
\hyp{H1} Coordination & \emph{Measured:} concentration of publishers. \emph{Not measured:} whether any datasets join across agencies; how agencies actually interact & 45 publishers, top two hold 59\% (partly an artefact of the search); only 6 datasets carry coordinates; cross-agency joinability was not tested\\
\hyp{H2} Procurement & \emph{Measured:} whether contracts or purchases appear in the estate; service formats; licence; readability. \emph{Not measured:} what was bought, or on what terms & No procurement dataset found; 2 datasets list an API format and none is a live feed; 17\% state no licence; 41\% of resources unreadable\\
\hyp{H3} Responsiveness & \emph{Measured:} age, spatial and time-indexed share of what the portal publishes. \emph{Not measured:} agencies' own operational systems & Median age 1.1\,y; every disaster dataset is a register or an archive, none is live\\
\hyp{H4} End-to-end & \emph{Measured:} link-by-link presence in the estate for two chains; shelter and tower records. \emph{Not measured:} warning delivery or outcomes & No drainage, evacuation-route or clean-air-room data found; 0 of 177 shelters carry a coordinate; 6 warning-tower sites\\\bottomrule
\end{tabularx}
\end{table*}

\paragraph{Data.} All inputs are public or were collected by the system; none required an agreement.
\begin{itemize}
\item \kw{Open-data estate.} Every dataset returned by a search for ``Chiang Mai'' on the national open-data portal on 19 September 2026: 311 datasets with 2{,}869 resources, each fetched and profiled; a resource is ``unreadable'' if the profiler could not parse it into a table (a PDF or image, a corrupt file, or a failed download; we did not separate the causes). A second catalogue of 682 resources (2 October) from the provincial portal was searched by title and tag. Absence statements below mean \emph{not found by title and tag search of these two sources}.
\item \kw{Disaster records.} The 11 datasets in the portal's disaster domain; the official temporary-shelter list (177 rows); the warning-tower list (12 rows).
\item \kw{System telemetry.} The system's own feed registry (32 entities, 4 November 2026); 2,576 camera frames received from three flood cameras (20 October to 4 November); 934 haze scores from one webcam, 917 paired with ground PM$_{2.5}$ (about four weeks); 104 level changes in the warning log (11 October to 3 November); a 271-commit build history (15 September to 3 November).
\item \kw{Burn seasons.} Monthly burned area, hotspot counts and PM$_{2.5}$ exceedance days for the 2024--25 and 2025--26 seasons, from the Hydro-Informatics Institute's burn-mapping data.
\end{itemize}
\paragraph{Shelters.} The official list has no coordinates, so we located each shelter with OpenStreetMap by name, accepting a point only if it lay inside the province outline and the stated district. Each point records its precision (building, subdistrict, district). We express sanitation as listed capacity per reported toilet and compare it with the Sphere medium-term planning ratio of 20 people per toilet (2018 handbook, printed p.\ 118 and Appendix 4, p.\ 146) \citep{sphere2018}. Sphere also permits an initial communal ratio of 50 in a rapid-onset crisis. Duration, accessibility, sex distribution, water supply and facility condition were not assessed, so the ratio is a scenario rather than a compliance test.
\paragraph{Reproducibility.} One script (\path{research/analysis/run_analysis.py}) regenerates the descriptive results and figures from archived files in the repository. Build history is frozen at commit \texttt{6075ee7}; milestone dates are commit-subject keyword matches, not release evidence. Counts, not significance tests, carry the argument: the estate is a census of one search, not a sample.

\section{The system}
\label{sec:system}
\begin{figure*}[!t]
\centering
\resizebox{\textwidth}{!}{\input{fig_chain.tikz}}
\caption{The chain from measurement to action. Boxes show the link; the grey line beneath names who publishes it, as far as we found, and whether a machine can read it. The red outline marks links we did not find in the estate.}\label{fig:chain}
\end{figure*}
We built \emph{CNX}, a public operations room for Chiang Mai (\url{https://cnx.nonarkara.org}), on five rules. CNX was not built from nothing: it began as a port of the second author\textquoteright s earlier provincial operations room for Lopburi (the first commit of this repository is titled a ``Lopburi-pattern port''), so the build times below measure an adaptation, not a first build.
\begin{enumerate}
\item \kw{Public sources only.} Read what agencies already publish. No memorandum with any agency. Some feeds need a free key or are subject to a provider\textquoteright s terms of use, which we did not negotiate and have not reviewed with counsel.
\item \kw{Say what you know and when.} Every number carries its source, its observation time and its failure mode. ``Unknown'' is never shown as ``calm''.
\item \kw{One verdict, thresholds from the source.} Rain bands follow the Thai Meteorological Department's 24-hour classes. A forecast may raise a level; it can never lower one that a measurement set.
\item \kw{Official first.} The meteorological service's own warnings lead the board. The board's own warnings leave as CAP messages marked \emph{restricted --- not an official public warning}.
\item \kw{Close the last link.} The board links hazard information to the official shelter register, with stated capacity and location precision; it does not verify opening status, safe access or admission.
\end{enumerate}
\noindent Figure~\ref{fig:chain} shows the chain the system tries to close. Its registry lists 32 entities from 29 distinct provider-label strings, not a count of independent agencies, including non-agency sources such as OpenStreetMap, news feeds and the board\textquoteright s own derived products (Figure~\ref{fig:registry}); 9 run through a relay on one machine, because their sources throttle or refuse the hosting platform's network or need work (such as image decoding) the platform cannot do.

\section{Results}
\label{sec:results}
\subsection{The estate: a record of the past, in files}
\begin{figure*}[!t]
\centering\includegraphics[width=\textwidth]{fig_estate}
\caption{The provincial open-data estate (n\,=\,311 datasets, 19 September 2026). (a) Years since each dataset was last updated. (b) Publishers by number of datasets. (c) Resources the profiler could not read.}\label{fig:estate}
\end{figure*}
\textbf{H1 --- coordination.} The estate has 45 publishers, but two of them (the Provincial Office and the Provincial Treasury) hold 59\% of datasets, and 20 publishers contribute exactly one (Figure~\ref{fig:estate}b; Herfindahl index 0.20). Because the search is for a province, provincial bodies dominate by construction, so concentration alone shows little. What matters for joining records is shared keys: only 6 datasets (1.9\%) carry a location, and we found no dataset about drainage, evacuation routes or clean-air rooms (Table~\ref{tab:links}). We did not test whether records from different agencies can in fact be joined.

\textbf{H2 --- procurement.} Across 311 datasets and 682 catalogue resources we found none about procurement, contracts or technology purchases. Two datasets list an API format (a provincial spending table and a provincial boundary layer), and none is a live feed. The estate is mostly files: 176 datasets offer CSV, 158 a spreadsheet and 44 a PDF; 1{,}182 of 2{,}869 resources (41\%) could not be parsed; 52 datasets (17\%) state no licence. This is an \emph{indirect} indicator. The portal does not publish what was bought, so it cannot show whether procurement favoured files over services, and a search for one province cannot show a national absence. It describes what a citizen or a neighbouring agency can reach, nothing more.

\textbf{H3 --- responsiveness.} The median dataset was 1.1 years old when baked; 131 of 311 had been updated within a year, and 37 were more than three years stale. All 11 datasets in the portal\textquoteright s disaster domain are registers or archives: injuries, deaths, damaged bridges, households affected, machinery and shelters (Table~\ref{tab:disaster}). The newest was updated on 25 July 2026, and the flood response plan is a PDF. An open-data portal is an archive by design, so this describes the portal, not the agencies\textquoteright\ operational systems, which we could not see.

\begin{table*}[!t]
\centering\small
\caption{The disaster domain of the estate (titles translated). ``Time'' = has a time field; ``Place'' = has coordinates.}\label{tab:disaster}
\begin{tabularx}{\textwidth}{@{}lXllcc@{}}
\toprule
\textbf{Dataset} & \textbf{Publisher} & \textbf{Updated} & \textbf{Format} & \textbf{Time} & \textbf{Place}\\\midrule
People injured in disasters & Chiang Mai Provincial Office & 2026-07-25 & CSV/XLSX & yes & no\\
Disaster-response personnel & Chiang Mai Provincial Office & 2026-07-25 & XLSX/CSV & yes & no\\
Warning-tower installation points & Chiang Mai Provincial Office & 2026-07-25 & CSV & yes & yes\\
Temporary shelters, Chiang Mai & Chiang Mai Provincial Office & 2026-07-25 & CSV/XLSX & no & no\\
Disaster-machinery register & Chiang Mai Provincial Office & 2026-07-25 & XLSX/CSV & yes & no\\
People killed in disasters & Chiang Mai Provincial Office & 2026-07-25 & CSV/XLSX & yes & no\\
Flood damage to bridges & Chiang Mai Provincial Office & 2026-07-25 & CSV/XLSX & yes & no\\
Households affected, by village & Chiang Mai Provincial Office & 2026-07-25 & CSV/XLSX & yes & no\\
Repeatedly burned areas & Lamphun Provincial Office & 2026-07-24 & XLSX/CSV & yes & no\\
Disaster-risk areas, 76 provinces & Dept.\ of Provincial Administration & 2026-03-24 & JSON & no & no\\
Provincial flood response plan & Dept.\ of Disaster Prevention and Mitigation & 2025-09-01 & PDF & no & no\\\bottomrule
\end{tabularx}
\end{table*}

\begin{table}[!t]
\centering\small
\caption{Datasets in the estate that match each link of the two chains (title and tag search, Thai and English terms). ``Library''  is the 311-dataset estate; ``Catalogue''  is the 682 provincial-portal resources.}\label{tab:links}
\begin{tabular}{@{}lcc@{}}
\toprule
\textbf{Link} & \textbf{Library} & \textbf{Catalogue}\\\midrule
Drainage network & 0 & 0\\
Evacuation routes & 0 & 0\\
Clean-air rooms & 0 & 0\\
Procurement / contracts & 0 & 0\\
API or live stream & 0 & 0\\
Warning towers & 1 & 1\\
Shelters & 1 & 2\\
Rainfall & 0 & 1\\
Flood events & 2 & 9\\
PM$_{2.5}$ / dust & 6 & 17\\
Fire and haze & 16 & 33\\\bottomrule
\end{tabular}
\end{table}

\subsection{Two chains, link by link}
Table~\ref{tab:chains} walks the two chains the research question names; Figure~\ref{fig:chain} draws the chain.
\begin{table*}[!t]
\centering\small
\caption{Rain-to-drainage and haze-to-shelter, link by link. ``Estate'' is what the open-data estate holds; ``CNX'' is what the system does with public sources.}\label{tab:chains}
\begin{tabularx}{\textwidth}{@{}p{2.1cm}XXX@{}}
\toprule
\textbf{Link} & \textbf{Rain $\to$ drainage} & \textbf{Haze $\to$ shelter} & \textbf{CNX}\\\midrule
Sense & Rain gauges published by the national water-data portal; machine-readable, but the portal throttles the hosting platform, so a relay fetches them & Ground sensors (university network, Pollution Control Department); satellite fire detections (375\,m) & Both read. Rain: 15-min relay. Fires: province-only counts\\
Forecast & Open weather-model grids; we found no official gridded forecast we could read & Model PM$_{2.5}$ and 850\,hPa transport & 48-h field and smoke trajectories, labelled as model\\
Assets & \textbf{None found in either catalogue} & \textbf{None found} & \emph{Not closed.} Named as a gap\\
Decide & Meteorological service publishes CAP alerts for provinces \citep{tmd2026cap} & We found no haze alert feed & Official alerts first; board verdict second\\
Push & CAP is available; no public push from CNX (it lacks the authority) & Same & CAP export, \emph{restricted}\\
Shelter & Official list with capacity and utilities, no coordinates & Same list; no air-filtration data & Shelters per district named in each warning; 45 located at the building\\
Verify & We found no published record of warnings scored against outcomes & Same & Warning log: 104 level changes since 11 Oct\\\bottomrule
\end{tabularx}
\end{table*}
\textbf{H4 --- end-to-end.} Of the seven links in each chain, the estate holds two in part: the shelter list (without coordinates) and the tower list (6 sites). We found no drainage network, evacuation routes or clean-air rooms. The meteorological service\textquoteright s alerts live outside the estate, in an open format. Where the system could join links from public sources, it did: gauges, fires, forecast and shelters sit in one verdict (Section~\ref{sec:system}). Where the link would be an asset register we could not find, it could not. Absence from a title-and-tag search of one province does not prove absence elsewhere --- a national drainage agency may hold such data --- so this is a finding about what a provincial user can reach.

\subsection{Shelters and warning towers}
\begin{figure}[!t]
\centering\includegraphics[width=\columnwidth]{fig_map}
\caption{Deduplicated shelter locations (174 records, by best available precision) and warning-tower sites in Chiang Mai. Other shelter summaries use 177 raw rows.}\label{fig:map}
\end{figure}
The official list has 177 rows in all 25 districts (three repeat an earlier row, so 174 are distinct), with a stated capacity of 98{,}775 people (median 200; the largest, 20{,}000, is one venue). \emph{None carries a coordinate.} After geocoding, 45 are located at the building, 77 at a subdistrict point and 55 at a district point (Figure~\ref{fig:shelters}b; Figure~\ref{fig:map}). For 75\% of rows, the derived point is only a district or subdistrict location, not a destination suitable for navigation. Even building-level matches require field verification.

\begin{figure}[!t]
\centering\includegraphics[width=\columnwidth]{fig_shelters}
\caption{(a) People per toilet at shelters that report any toilets; the line is the reference value of 20. (b) Best location we can state for each shelter.}\label{fig:shelters}
\end{figure}
Sanitation is thinner than capacity. 31 shelters (18\%) report zero toilets; of the 146 that report any, 73 (50\%) exceed 20 people per toilet (median 22.5). If each shelter hosted only as many people as its toilets serve at that reference value, the stated 98{,}775 places would fall to 40{,}739 (41\%). This is a planning scenario, not an estimate of usable capacity or a mathematical lower bound: a reported zero may indicate missing data, and facilities may be shared with neighbouring buildings. The capacity sum also includes three duplicate rows.

The warning-tower list is an asset register kept by the National Disaster Warning Center for towers that alert the public to tsunami, flood and landslide. It has 12 rows, which are the same 6 sites listed in each of two yearly registers (2024 and 2025); we therefore count 6 sites. Three of the six stand within a few kilometres of one another in the city. The nearest tower to any of the 45 building-level shelters is a median 25\,km away; 3 lie within 5\,km. The data do not say how far a tower\textquoteright s sound carries, and the list may be partial, so these distances are descriptive, not a coverage estimate.

\subsection{What the system took to build}
\begin{figure}[!t]
\centering\includegraphics[width=\columnwidth]{fig_build}
\caption{Cumulative commits across the build. Markers indicate the first commit-subject keyword matches for rain, forecast and official warnings. These are development-history proxies, not verified deployment dates.}\label{fig:build}
\end{figure}
The repository\textquoteright s first commit was on 15 September 2026, and it began as a port (Section~\ref{sec:system}). By 3 November 2026 it had 271 commits on 36 active days within a seven-week calendar span. The first commit subjects matching rain, forecast and CAP 1.2 occur on days 4, 16 and 26 respectively (Figure~\ref{fig:build}). Keyword matches may predate functioning features and do not establish when users received them. No agency was asked for anything; this is evidence of feasibility for one adapted system, not of the cost of a first build.

\begin{figure}[!t]
\centering\includegraphics[width=\columnwidth]{fig_registry}
\caption{The feed registry: cadence against the staleness limit each feed declares. Markers show how the feed reaches the system; colours show its kind.}\label{fig:registry}
\end{figure}
\textbf{Honest weaknesses.} Of 32 entities, only 1 is ground-truthed and 5 cross-checked; 26 (81\%) have no validation at all. The system shows data; it does not certify it. 9 entities depend on one relay machine (28\%), a single point of failure the board reports as ``unavailable'', never as ``calm''.

\subsection{Operating record}
In the archived 4 November 2026 registry snapshot, 24 of 25 feeds with a declared limit were inside it; the one outside was a satellite product with no scene published for days, shown as unavailable. For each flood camera, distinct-frame counts divided by the number of nominal 15-minute intervals between its first and last recorded observations were 41--72\% over approximately 14.7 days (Figure~\ref{fig:cameras}); the longest silence was 36 hours. Frames arrived a median of 18--21 minutes after capture.
\begin{figure}[!t]
\centering\includegraphics[width=\columnwidth]{fig_cameras}
\caption{Frames received from three flood cameras, 20 October--4 November 2026. Gaps indicate no distinct frame in the archived series; source, transport, relay and unchanged-image causes are not isolated.}\label{fig:cameras}
\end{figure}
One webcam\textquoteright s haze scoring has been paired with 917 associated ground PM$_{2.5}$ readings. The readings ranged 1--21\,\textmu g/m$^3$ (median 6); 20 exceeded 15 and none exceeded 37.5 (the WHO\textquoteright s 24-hour guideline and a 24-hour interim target; we compare single readings with a 24-hour value, so these counts do not measure guideline exceedance) \citep{who2021aqg}. This PM$_{2.5}$ range contains no reading above 37.5. It does not establish the absence of a haze event. Pairing tolerance, spatial representativeness, independent labels and a decision threshold were not evaluated here; sensitivity, specificity and false-alarm rates therefore remain unknown. Repeated observations from one webcam are not independent validation samples.

The board's warning log recorded 104 level changes (83 forecast-PM$_{2.5}$ district warnings and 21 province-verdict changes), each with its reason, source and build. It is the beginning of the record that Table~\ref{tab:chains} found missing: warnings that can later be scored against outcomes.

\subsection{Burn seasons: the pattern the chain should catch}
\begin{figure}[!t]
\centering\includegraphics[width=\columnwidth]{fig_burn}
\caption{Monthly burned area (farm and forest) in two seasons.}\label{fig:burn}
\end{figure}
We hold two seasons, each covering November to April; hotspot counts are missing for November and December in both. In 2025--26, March and April carried 45 of the 47 PM$_{2.5}$-exceedance days; in 2024--25, 37 of 49. Burned area was lower in 2025--26 (413 against 527 thousand rai) while hotspot detections were higher (2{,}072 against 673, over different months). The two series measure different things, and the data owner warns that monthly area counts a field burned twice in both months, so they are not expected to reconcile. We show them to make one point: the two archived seasons concentrate exceedance days in March and April; two seasons do not establish predictive skill or source attribution. A maintained monitoring service could be evaluated against this seasonal pattern.

\section{Discussion}
\paragraph{What integration demonstrates.} The prototype reduced some practical barriers to assembling a shared view; it did not demonstrate that institutional gaps were bypassed. Coordination: the system read each provider\textquoteright s public output and joined it itself, without waiting for an agreement. Procurement: it made no purchase, so it shows that a prototype can avoid the procurement gap, not that procurement is the cause; the system still runs on commercial hosting and free-tier feeds. Responsiveness: 25 of its 32 entities declare a staleness limit and show their age. The end-to-end chain remains unverified: display and CAP export do not prove delivery, acknowledgement, evacuation or a safe shelter arrival.

\paragraph{What cannot be bypassed.} Three things stayed out of reach, and they are the honest limit of this work. (1) \emph{Authority.} A public operations room can advise, and the board says plainly that it issues no evacuation order and no public alert. The national system for such orders is set by the Disaster Prevention and Mitigation Act \citep{thailand2007dpmact}; we claim no part of it. (2) \emph{Assets.} Nobody can route water to a drain that nobody has mapped. (3) \emph{Validation and staffing.} Most feeds are unvalidated against ground truth, and a screen with no named duty officer is a screen. The system's own manual still lists its owner, operator and duty officer as ``to be assigned''.

\paragraph{A mechanism for future study.} One possible explanation is that project-based publication favours static deliverables over maintained services. Static records alone do not establish who receives a warning or acts on it. Whether procurement and organisational incentives explain this pattern requires contract evidence and interviews. This is a hypothesis for investigation alongside the literature on public IT \citep{goldfinch2007,heeks2006,pahlka2023}, not a result.

\paragraph{Policy implications.} Four policy priorities follow; their implementation costs were not estimated. Publish three records the chain needs: shelter coordinates, drainage geometry, and tower locations with audible range. Require, in each technology purchase, a documented machine interface and a licence \citep{edler2007,mergel2019}. Publish a freshness limit for every public feed, so that silence is visible. Carry every official warning in CAP \citep{oasis2010cap}, as the meteorological service\textquoteright s feed already does \citep{tmd2026cap}.

\paragraph{Threats to validity.}\label{sec:limits} (1) \emph{Design.} Single case, single province, one adapted system built by the second author. The audit is descriptive; H1--H3 are not tested, and H1\textquoteright s cross-agency joins and H2\textquoteright s procurement are not measured at all. A test would need interviews and procurement records. (2) \emph{Sampling.} The estate is the result of one search; absence means absence from that search, and provincial publishers dominate it by construction. (3) \emph{Evaluation.} The system is evaluated for implementation feasibility, a feed-freshness snapshot and archived frame receipt only; application uptime was not measured. We have no user study, no lead-time or accuracy test of its warnings, one webcam, and no independently labelled haze events; 26 of 32 feeds have no validation. (4) \emph{Data.} Geocoded shelter points rest on OpenStreetMap coverage; the tower list may be partial; camera figures count distinct frames, so a frozen image counts as silence; the portal data may have changed since 19 September. (5) \emph{Interests.} The second author works for the agency that promotes the programme discussed and built the system; the first is affiliated with the case study\textquoteright s sponsor (Declarations). We report gaps in the second author\textquoteright s own sector for that reason. We show the gaps are present, measurable and partly bridgeable; we do not show that they \emph{cause} the stalled programmes.

\section{Conclusion}
In Chiang Mai, public sources supported an adapted operations room joining measurements, official warnings and shelter records within a seven-week repository history. The descriptive audit identifies missing coordinates, absent records within the searched catalogues, and incomplete validation. It does not demonstrate an operationally closed warning-to-shelter chain, improved safety outcomes or the causes of national programme performance. The next evaluation needs named duty officers, verified shelter access, message-delivery and acknowledgement logs, and independently labelled flood and haze events. Publishing and maintaining the missing asset records is a policy and operational task alongside the software work.

\section*{Declarations}
\small
\textbf{Competing interests.} P.T.\ is affiliated with RCAD, which sponsors the case study, and with the Faculty of Engineering, Chiang Mai University. N.A.\ is employed by the Digital Economy Promotion Agency (depa), which promotes the Smart City Thailand programme discussed here, and designed and built the CNX system. Whether either organisation reviewed this manuscript, and any funding for the work, are \textbf{to be stated by the authors before submission}. The views are the authors\textquoteright\ own and not those of any organisation.\\
\textbf{Data and code availability.} Data, analysis script and paper source: \path{research/} in the CNX repository at commit \texttt{bad3fce}, and \url{https://cnx.nonarkara.org/cnx/research}. An archived copy with a DOI is \textbf{to be deposited} before submission. The portal data were baked on 19 September and 2 October 2026 and may have changed.\\
\textbf{AI assistance.} Claude (Anthropic) assisted the original analysis, figures and drafting; Codex (OpenAI) assisted this revision. Exact model versions and usage details require author confirmation. The authors remain responsible for the content. AI-assisted editorial and numerical checks are not independent scientific validation. References were checked against Crossref, publisher or official pages; the bibliography records differing levels of verification. Metadata checks do not establish that a source supports an argument. Authors must check the cited passages before submission (see \path{references-verified.bib}).\\
\textbf{Author contributions.} \textbf{To be confirmed by the authors before submission.}\\
\textbf{Ethics and terms of use.} The released research series contain derived timestamps, image scores and associated PM$_{2.5}$ values, rather than camera images or public-post text. This is not a privacy audit of the operational system: public posts, links and camera imagery may contain identifiable information. Lawful basis, retention, access controls and any ethics-review requirement remain to be assessed before expanding collection. The Personal Data Protection Act is relevant to that assessment \citep{thailand2019pdpa}. The relay fetches public endpoints at a modest rate and does not bypass authentication; a formal terms-of-use review of each provider has not been done.

{\small\bibliographystyle{abbrvnat}\bibliography{refs}}
\end{document}
